\subsection*{Problema 3 (ejercicio 4)}
\emph{Solución.}
\begin{enumerate}
\item \(\left\{n\atop2\right\}=2^{n-1}-1\). Elegir para cada \(i>1\) si va con el 1, menos el caso en que todos lo hacen.
\item \(\left\{n\atop n-1\right\}=\binom n2=\frac{n(n-1)}2\). Un bloque de 2 y \(n-2\) simples.
\item \(\left\{n\atop n-2\right\}=\binom n3 + 3\binom n4\). Ya que hay
  \begin{itemize}
  \item Un bloque de 3: \(\binom n3\).
  \item Dos bloques de 2: elegimos 4 elementos \(\binom n4\) y hay 3 formas de emparejarlos.
  \end{itemize}
\end{enumerate}
